%%%%%%%%%%%%%%%%%%%%%%%%%%%%%%%%%%%%%%%%%%%%%%%%%%%%%%%%%%%%
% (01) Input Settings
%%%%%%%%%%%%%%%%%%%%%%%%%%%%%%%%%%%%%%%%%%%%%%%%%%%%%%%%%%%%
\input{settings.tex}

%%%%%%%%%%%%%%%%%%%%%%%%%%%%%%%%%%%%%%%%%%%%%%%%%%%%%%%%%%%%
% (02) DOCUMENT INFORMATION - FROM YAML METADATA
%%%%%%%%%%%%%%%%%%%%%%%%%%%%%%%%%%%%%%%%%%%%%%%%%%%%%%%%%%%%
\newcommand{\DocTitle}{مقدمه‌ای بر شی‌گرایی}

\newcommand{\CourseTitle}{برنامه‌سازی پیشرفته و کارگاه}

\newcommand{\Semester}{بهار ۱۴۰۵}

\newcommand{\FirstProfessor}{دکتر احمدیان}

\newcommand{\SecondProfessor}{}

\newcommand{\Writers}{شهاب گریوانی، یونس کاظمی، آرمان حسینی، نگین موسی
یی جو، نوا جغتایی}

\newcommand{\Faculty}{دانشکده ریاضی و علوم کامپیوتر}

\newcommand{\University}{دانشگاه صنعتی امیرکبیر}

%%%%%%%%%%%%%%%%%%%%%%%%%%%%%%%%%%%%%%%%%%%%%%%%%%%%%%%%%%%%
% DOCUMENT BODY
%%%%%%%%%%%%%%%%%%%%%%%%%%%%%%%%%%%%%%%%%%%%%%%%%%%%%%%%%%%%
\begin{document}
	
%%%%%%%%%%%%%%%%%%%%%%%%%%%%%%%%%%%%%%%%%%%%%%%%%%%%%%%%%%%%
% === TITLE PAGE ===
%%%%%%%%%%%%%%%%%%%%%%%%%%%%%%%%%%%%%%%%%%%%%%%%%%%%%%%%%%%%
\begin{titlepage}
	\centering
	\vspace*{1cm}
	
	%— Smaller logo (20% of text width)
	\begin{figure}[H]
		\centering
		
\includegraphics[width=0.20\textwidth]{images/logo-fa.png}
	\end{figure}
	
	\vspace{1cm}
	
	{\fontsize{18}{28}\selectfont \textbf{\CourseTitle} }
	
	\vspace{3cm}
	
	{\fontsize{22}{28}\selectfont \textbf{\DocTitle} }
	
	\vspace{3cm}
	
	{\fontsize{14}{22}\selectfont \rl{نگارش:}}\\
	\vspace{0.5cm}
	{\fontsize{16}{22}\selectfont \textbf{ \Writers } }

	
	\vspace{1cm}
	
	{\fontsize{14}{22}\selectfont \rl{استاد درس:}}\\
	\vspace{0.5cm}
	{\fontsize{16}{22}\selectfont \textbf{ \FirstProfessor } }
	
	
	\vfill
	
	{\fontsize{14}{18}\selectfont \textbf{\Faculty} }\\
	\vspace{0.5cm}
	{\fontsize{14}{18}\selectfont \textbf{\University} }
	\vspace{0.5cm}
	
	\par
	
	{\fontsize{14}{18}\selectfont \textbf{\Semester} }
	
\end{titlepage}

%— Switch to our fancy header/footer from this page onward
\clearpage

%%%%%%%%%%%%%%%%%%%%%%%%%%%%%%%%%%%%%%%%%%%%%%%%%%%%%%%%%%%%
% === TABLE OF CONTENTS ===
%%%%%%%%%%%%%%%%%%%%%%%%%%%%%%%%%%%%%%%%%%%%%%%%%%%%%%%%%%%%
\tableofcontents
\clearpage
\pagestyle{aftertitle}

%%%%%%%%%%%%%%%%%%%%%%%%%%%%%%%%%%%%%%%%%%%%%%%%%%%%%%%%%%%%
% === MAIN CONTENT ===
%%%%%%%%%%%%%%%%%%%%%%%%%%%%%%%%%%%%%%%%%%%%%%%%%%%%%%%%%%%%
\chapter{\DocTitle}
\section{مقدمه}

در این داکیومنت و داکیومنت‌های بعدی، با شی‌گرایی آشنا می‌شویم. شی‌گرایی یا
\lr{object-oriented programming} (به اختصار \lr{OOP} )، طرز فکری است
برای تقسیم یک برنامهٔ بزرگ، به واحدهای کوچک‌تر به اسم \lr{class .} خیلی از
چیزهایی که قبلا بهتون گفته بودیم «در آینده برایتان توضیح می‌دهیم» را قرار
است این‌جا یاد بگیرید!

پیشنهاد می‌کنیم در طول مطالعه‌ی این داکیومنت با تدریس یارها در ارتباط
باشید و اگر سوالی داشتید حتما از آن‌ها بپرسید، زیرا این یکی از مهم‌ترین
داکیومنت‌های شماست و پایه‌ایست برای اکثر مطالبی که در ادامه‌ی ترم باهم مورد
بحث قرار می‌دهیم مانند ارث‌بری، کپسوله‌سازی و مباحث مشابه .

\section{تایپ‌ها در جاوا}

تا به اینجا با بسیاری از \lr{type} ها مثل \lr{\texttt{int}} ،
\lr{\texttt{float}} ، \lr{\texttt{String}} ، \lr{\texttt{JFrame}} ،
\lr{\texttt{ArrayList}} و \ldots{} کار کرده‌اید و با آن‌ها آشنا شده‌اید؛ هر
برنامه‌نویسی در جاوا، می‌تواند به راحتی تایپ‌های جدید برای خودش بسازد و به
همین خاطر، احتمالا تا به حال میلیون‌ها تایپ در جاوا نوشته شده‌اند. شما نیز
در این داک قرار است یاد بگیرید که چطور می‌توانید تایپ‌های جدید بسازید.

همان‌طور که در جلسه‌ی صفر دیدید، تایپ‌های \lr{\texttt{int}} ،
\lr{\texttt{float}} ، \lr{\texttt{double}} ، \lr{\texttt{long}} و
\lr{\texttt{char}} ، به اسم «تایپ‌های اولیه» یا \lr{primitive types}
شناخته می‌شوند. فهرست کامل آن‌ها به شکل زیر است:

\begin{table}[H]
\centering
\setLTR
\begin{tabular}{lll}
\hline
\textbf{\lr{Type}} & \textbf{\lr{Definition}} & \textbf{\lr{Approximate range or precision}} \\
\hline
\texttt{boolean} & \lr{Logical Value} & \texttt{true} \lr{or} \texttt{false} \\
\texttt{char} & \(16\)\lr{-bit, Unicode character} & \(64k\) \lr{characters} \\
\texttt{byte} & \(8\)\lr{-bit, signed integer} & \(-128\) \lr{to} \(127\) \\
\texttt{short} & \(16\)\lr{-bit, signed integer} & \(-32, 768\) \lr{to} \(32,676\) \\
\texttt{int} & \(32\)\lr{-bit, signed integer} & \(-2.1\mathrm{e}9\) \lr{to} \(2.1\mathrm{e}9\) \\
\texttt{long} & \(64\)\lr{-bit, signed integer} & \(-9.2\mathrm{e}18\) \lr{to} \(9.2\mathrm{e}18\) \\
\texttt{float} & \(32\)\lr{-bit, IEEE 754, floating-point value} & \(6-7\) \lr{significant decimal places} \\
\texttt{double} & \(64\)\lr{-bit, IEEE} 754 & \(15\) \lr{significant decimal places} \\
\hline
\end{tabular}
\end{table}

این تایپ‌ها، مانند تیکه‌های کوچک لگو هستند که تایپ‌های پیچیده‌تر مانند
\lr{\texttt{String}} ، \lr{\texttt{JFrame}} ، \lr{\texttt{ArrayList}} و
هر تایپ دیگری را تشکیل می‌دهند. این تایپ‌های پیچیده‌تر، \lr{reference type}
نام دارند. در ادامه یک مثال از آن‌ها را مورد بررسی قرار می‌دهیم. کد زیر را
در \lr{IntelliJ} کپی کنید:

\begin{codebox}{java}
public class Main {
    public static void main(String[] args) {  
        String str = "Hello World";  
        System.out.println(str);  
    }  
}
\end{codebox}

موس را روی \lr{\texttt{String}} ببرید و کلیک راست کنید. از صفحه‌ای که باز
می‌شود، منوی \lr{Go To} و بعد از آن، \lr{Declaration or Usages} را انتخاب
کنید:

\begin{figure}[H]
\centering
\includegraphics[width=0.8\textwidth]{images/oop-1-1.png}
\end{figure}

صفحه‌ی جدیدی برایتان باز می‌شود که شامل پیاده‌سازی تایپ
\lr{\texttt{String}}است:

\begin{figure}[H]
\centering
\includegraphics[width=0.8\textwidth]{images/oop-1-2.png}
\end{figure}

احتمالا شما هنوز این کد را کامل نمی‌فهمید، ولی در این کد نشان داده شده
است که تایپ \lr{\texttt{String}} از کنار هم گذاشتن چه تایپ‌های دیگری درست
شده است. مثلا در خط ۱۶۰ همین تصویر می‌بینید که \lr{\texttt{String}} در
خودش یک آرایه از \lr{\texttt{byte}} ها دارد که اسمش \lr{\texttt{value}}
است. پایین‌تر می‌بینید که در خودش، یک \lr{\texttt{int}} به نام
\lr{\texttt{hash}} دارد. اگر دوست دارید، کمی این کد را بررسی کنید. با آن
که تیکه‌هایی از آن برایتان غریبه هستند، بخش‌هایی هم برای شما آشنا به نظر
می‌آیند.

\section{کلاس‌ها}

کلاس‌ها، به شما اجازه می‌دهند تا \lr{\texttt{type}} های جدیدی درست کنید.
برای شروع، با هم یه کلاس جدید به اسم \lr{\texttt{Student}} درست می‌کنیم.
کد زیر را در \lr{IntelliJ} بنویسید:

\begin{codebox}{java}
public class Main {
    public static void main(String[] args) {  
        System.out.println("Hello World!");  
    }  
}

class Student {
    public int age; 
    public String name;
    public String studentID;  
}
\end{codebox}

در این کد، ما یک \lr{\texttt{class}} جدید، به نام \lr{\texttt{Student}}
تعریف کرده‌ایم. این کلاس، از چهار متغیر تشکیل شده است:

\begin{itemize}
\tightlist
\item
  متغیر \lr{\texttt{age}} : متغیری از جنس \lr{\texttt{int}} که سن دانشجو
  را نشان می‌دهد.
\item
  متغیر \lr{\texttt{name}} : متغیرِ \lr{\texttt{String}} ای که اسم دانشجو
  را نشان می‌دهد.
\item
  متغیر \lr{\texttt{studentID}} : متغیر \lr{\texttt{String}} ای که
  شماره‌ی دانشجویی شخص را نشان می‌دهد.
\end{itemize}

متغیرهایی که یک کلاس را تشکیل می‌دهند، \lr{\texttt{field}} های آن کلاس
نام دارند. نگران کلیدواژه‌ی \lr{\texttt{public}} نباشید، در جلسات آینده
مفهوم آن را توضیح خواهیم داد. برای الان، لازم است که آن را پیش از همه‌ی
\lr{\texttt{field}} های کلاس‌هایتان قرار دهید. حالا با کلیدواژهٔ
\lr{\texttt{new}} ، یک دانشجوی جدید به اسم قلی بسازید و فیلدهای آن را
مقداردهی کنید:

\begin{codebox}{java}
public static void main(String[] args) {  
    Student gholi = new Student();  
  
    gholi.age = 20;  
    gholi.name = "Gholi";  
    gholi.studentID = "40413099";  
}
\end{codebox}

سپس مشخصات این دانشجو را چاپ کنید:

\begin{codebox}{java}
System.out.println("New student:");  
System.out.println("\t+ Name: " + gholi.name);  
System.out.println("\t+ StudentID: " + gholi.studentID);  
System.out.println("\t+ Age: " + gholi.age);
\end{codebox}

کد خود را اجرا کنید. خروجی زیر را می‌بینید:

\begin{codebox}{text}
New student:  
    + Name: Gholi  
    + StudentID: 40413099  
    + Age: 20
\end{codebox}

می‌بینید که مشخصات قلی، به درستی روی صفحه چاپ می‌شود. تبریک! شما اولین
\lr{type} خود را ساختید و یک متغیر از جنس آن درست کردید. می‌خواهیم به این
\lr{type} جدید، چیزهای بیشتری اضافه کنیم. به \lr{\texttt{Student}} فیلدی
به اسم \lr{\texttt{grades}} از جنس \lr{\texttt{ArrayList<Double>}} اضافه
کنید. این فیلد، نمرات دانشجو در درس‌های مختلف را نشان می‌دهد:

\begin{codebox}{java}
public ArrayList<Double> grades;
\end{codebox}

حالا، وقتی دارین فیلدهای مختلف قلی رو توی \lr{\texttt{main}} مقداردهی
می‌کنید، این آرایه هم با یک آرایهٔ خالی مقداردهی کنید:

\begin{codebox}{java}
gholi.age = 20;  
gholi.name = "Gholi";  
gholi.studentID = "40413099";  
gholi.grades = new ArrayList<Double>();
\end{codebox}

بعد از این کار، چند نمره‌ی رندوم به قلی بدهید:

\begin{codebox}{java}
gholi.grades.add(20.0);  
gholi.grades.add(17.0);  
gholi.grades.add(18.0);  
gholi.grades.add(0.0);
\end{codebox}

حالا، جایی که دارید مشخصات قلی را چاپ می‌کنید، نمرات را هم چاپ کنید:

\begin{codebox}{java}
System.out.println("New student:");  
System.out.println("\t+ Name: " + gholi.name);  
System.out.println("\t+ StudentID: " + gholi.studentID);  
System.out.println("\t+ Age: " + gholi.age);  
System.out.print("\t+ Grades: ");
for (var grade: gholi.grades) {  
    System.out.print(grade + ", ");  
}
\end{codebox}

کدتان را اجرا کنید. خروجی شما باید به این شکل باشد:

\begin{codebox}{text}
New student:  
    + Name: Gholi  
    + StudentID: 40413099  
    + Age: 20  
    + Grades: 20.0, 17.0, 18.0, 0.0,
\end{codebox}

می‌خواهیم به \lr{\texttt{Student}} ، قابلیت محاسبه‌ی معدل هم بدهیم. برای
این کار، متدی جدید به اسم \lr{\texttt{getAverageGrade}} در
\lr{\texttt{Student}} تعریف می‌کنیم:

\begin{codebox}{java}
class Student {
    public int age;
    public String name;
    public String studentID;
    public ArrayList<Double> grades;

    public double getAverageGrade() { 
        if (grades.size() == 0) {
            return 0;  
        }

        double gradeSum = 0;
        for (double grade: grades) {  
            gradeSum += grade;  
        }

        return gradeSum / grades.size();  
    }  
}
\end{codebox}

جایی از کدتان که در آن مشخصات \lr{\texttt{gholi}} را چاپ می‌کنید، خطوط
زیر را اضافه کنید:

\begin{codebox}{java}
System.out.println();  
System.out.println("\t+ Average Grade: " + gholi.getAverageGrade());
\end{codebox}

کدتان را دوباره اجرا کنید. خروجی‌ای مثل زیر می‌بینید:

\begin{codebox}{text}
New student:  
    + Name: Gholi  
    + StudentID: 40413099  
    + Age: 20  
    + Grades: 20.0, 17.0, 18.0, 0.0,   
    + Average Grade: 13.75
\end{codebox}

لحظه‌ای به خود تابع \lr{\texttt{main}} توجه کنید. می‌بینید که خود آن نیز
در کلاسی به اسم \lr{\texttt{Main}}‍است! شما در تمام این مدت داشتید کلاس
\lr{\texttt{Main}} را تعریف می‌کردید، بدون آن که خبر داشته باشید:

\begin{codebox}{java}
public class Main {
    public static void main(String[] args) {
        // Your code here
    }  
}
\end{codebox}

شما حتی می‌توانید از این کلاس هم یک متغیر درست کنید:

\begin{codebox}{java}
public class Main {
    public static void main(String[] args) {  
        Main a = new Main();  
    }  
}
\end{codebox}

البته که متغیر ساختن از جنس \lr{\texttt{Main}} خیلی کار خوبی نیست! ولی
جالب است که تا همین الآن هم شما از \lr{\texttt{class}} ها استفاده
می‌کردید، بدون آن که بدانید. =)))

در ادامه، کمی رسمی‌تر و قدم به قدم‌تر کلاس‌ها را بررسی می کنیم.

\subsection{تعریف کلاس}

تعریف کلاس، کار سختی نیست، فقط از کلیدواژهٔ \lr{\texttt{class}} استفاده
کنید و اسم کلاستان را بنویسید:

\begin{codebox}{java}
class Student {
    // Everything an student can do
}
\end{codebox}

هر چیزی که بین دو براکت قرار گیرد، متعلق به آن کلاس است. کلاس‌ها، به شما
اجازه می‌دهند که کدهای خیلی خیلی بزرگ را به تیکه‌های کوچیک‌تر تقسیم کنید و
با این کار، برنامه‌های منظم‌تر و بهتری داشته باشید.

یادآوری می‌کنیم که \lr{type} هایی که به وسیلهٔ \lr{\texttt{class}} ها
تعریف می‌شوند، همگی از \lr{reference type} ها محسوب می‌شوند.

\section{ایجاد \lr{Object}}

به متغیرهایی که از جنس یک \lr{reference type} هستند، \lr{object} یا
\lr{instance} می‌گویند.در کد زیر، قلی و ممد و سلطان همگی \lr{object} هایی
از جنس \lr{\texttt{Student}} ان:

\begin{codebox}{java}
public static void main(String[] args) {  
    Student gholi;  
    Student mamad;  
    Student soltan;  
}
\end{codebox}

تلاش کنید این کد را اجرا کنید. می‌بینید که طبیعتا بدون هیچ مشکلی اجرا
می‌شود. شما توی این کد سه \lr{object} جدید تعریف کردید و هیچ‌جایی آن‌ها را
مقداردهی نکردید و جایی هم از آن‌ها استفاده نشده است.

برای این که آن‌ها را مقداردهی کنید، از کلید‌واژهٔ \lr{new} استفاده کنید:

\begin{codebox}{java}
Student gholi = new Student();  
Student mamad = new Student();  
Student soltan = new Student();
\end{codebox}

حالا فیلدهای ان‌ها را مقداردهی کنید:

\begin{codebox}{java}
gholi.age = 21;  
gholi.name = "Gholi";  
gholi.studentID = "40513089";  
gholi.grades = new ArrayList<>();  
  
mamad.age = 25;  
mamad.name = "Mamad";  
mamad.studentID = "40513090";  
mamad.grades = new ArrayList<>();  
  
soltan.age = 103;  
soltan.name = "Soltan";  
soltan.studentID = "40513091";  
soltan.grades = new ArrayList<>();
\end{codebox}

همانطور که می‌بینید ما در این کد به اجبار کد های تکراری زیادی زده‌ایم. کدی
که مربوط به بخش مقداردهی \lr{field} های یک \lr{object} است سه بار تکرار
می‌شود پس بهتر است این کدها را به یک متد تبدیل کنیم.
\lr{\texttt{Constructor}} ها، متدهایی هستند که دقیقا برای همین کار، به
کمک ما می‌آیند.

\subsection{\lr{Constructor} ها}

به کلاس \lr{Student} ، متد زیر را اضافه کنید:

\begin{codebox}{java}
public Student(int newStudentAge, String newStudentName, String newStudentID) {
    age = newStudentAge;
    name = newStudentName;
    studentID = newStudentID;

    grades = new ArrayList<>();  
}
\end{codebox}

می‌بینید که ظاهر این متد، کمی متفاوت است! نوع خروجی‌اش مشخص نشده و جایی از
کلیدواژه‌ی \lr{\texttt{return}} استفاده نکرده است. به این متد،
\lr{\texttt{Constructor}} می‌گویند و وظیفه‌ی آن، مقداردهی یک \lr{object}
جدید از کلاس \lr{\texttt{Student}} ‍است. در \lr{constructor} بالا، ما
برای ساخت یه \lr{\texttt{Student}} جدید، سن، اسم و شماره دانشجویی‌اش را
ورودی گرفته‌ایم و با استفاده از آن‌ها فیلدهای \lr{\texttt{Student}} را
برای \lr{object} جدیدمان مقداردهی کرده‌ایم.

برای این که از این \lr{\texttt{constructor}} استفاده کنید، کد توی
\lr{\texttt{main}} را با کد زیر جایگرین کنید:

\begin{codebox}{java}
public static void main(String[] args) {  
    Student gholi = new Student(21, "Gholi", "40513089");  
    Student mamad = new Student(25, "Mamad", "40513090");  
    Student soltan = new Student(103, "Soltan", "40513091");  
}
\end{codebox}

می‌بینید که با استفاده از یه \lr{\texttt{constructor}} خوب، چقدر کد
\lr{\texttt{main}} کوتاه‌تر و تمیزتر شد! حالا متد زیر را به
\lr{\texttt{Student}} اضافه کنید:

\begin{codebox}{java}
public void printInfo() {  
    System.out.println("Student info:");  
    System.out.println("\t+ Name: " + name);  
    System.out.println("\t+ StudentID: " + studentID);  
    System.out.println("\t+ Age: " + age);  
    System.out.print("\t+ Grades: ");
    for (var grade: grades) {  
        System.out.print(grade + ", ");  
    }  
    System.out.println();  
    System.out.println("\t+ Average Grade: " + getAverageGrade());  
}
\end{codebox}

و توی \lr{\texttt{main}} ، این متد را صدا بزنید تا اطلاعات ممد و قلی و
سلطان چاپ شود:

\begin{codebox}{java}
gholi.printInfo();  
mamad.printInfo();  
soltan.printInfo();
\end{codebox}

خروجی‌ای مثل زیر می‌بینید:

\begin{codebox}{java}
Student info:  
    + Name: Gholi  
    + StudentID: 40513089  
    + Age: 21  
    + Grades:   
    + Average Grade: 0.0  
Student info:  
    + Name: Mamad  
    + StudentID: 40513090  
    + Age: 25  
    + Grades:   
    + Average Grade: 0.0  
Student info:  
    + Name: Soltan  
    + StudentID: 40513091  
    + Age: 103  
    + Grades:   
    + Average Grade: 0.0
\end{codebox}

می‌بینیم که \lr{constructor} ها کار کردند و کد ما هم، کوتاه و تمیز شد!

\subsection{\lr{Default Constructor}}

پیش از آن که ما برای کلاس \lr{\texttt{Student}} کانستراکتور تعریف کنیم،
می‌توانستیم با استفاده از کلیدواژهٔ \lr{\texttt{new}} از آن \lr{object}
بسازیم:

\begin{codebox}{java}
Student gholi = new Student();
\end{codebox}

خب، کمی بهش فکر کنید. چگونه این اتفاق افتاد؟ ما
که\lr{\texttt{constructor}} ای برای \lr{\texttt{Student}} نداشتیم! با
نوشتن این کد، دقیقا چه \lr{\texttt{constructor}} ای صدا زده می‌شود؟

وقتی یک کلاس \lr{\texttt{constructor}} نداشته باشد، جاوا به طور پیش‌فرض
یک \lr{\texttt{constructor}} خالی برایش می‌نویسد. این
\lr{\texttt{constructor}} به شکل زیر است:

\begin{codebox}{java}
public Student() {}
\end{codebox}

در این کانستراکتور هیچ اتفاقی نمی‌افتد و هیچ \lr{field} ای مقداردهی
نمی‌شود. صرفاً وجود دارد تا شما بتوانید به راحتی با استفاده از
\lr{\texttt{new Student()}} ، آبجکت‌های جدیدی از این کلاس بسازید. به محض
آن که شما اولین \lr{\texttt{constructor}} کلاس \lr{\texttt{Student}} را
بنویسید، دیگر نمی‌توانید از \lr{default constructor} جاوا استفاده کنید.
در کد \lr{main} فعلیتان این خط را بنویسید و آن را اجرا کنید:

\begin{codebox}{java}
Student pedram = new Student();
\end{codebox}

خطایی مانند زیر می‌گیرید:

\begin{figure}[H]
\centering
\includegraphics[width=0.8\textwidth]{images/oop-1-3.png}
\end{figure}

از آن‌جایی که کلاس \lr{\texttt{student}} ، الآن یک
\lr{\texttt{constructor}} دارد، جاوا برای شما کانستراکتور دیفالت را درست
نمی‌کند.

\subsection{کلاس‌ها با چند کانستراکتور}

شما می‌توانید برای رفع مشکلی که بالاتر به آن برخورد کردید، خودتان یک
کانستراکتور خالی برای \lr{\texttt{Student}} بنویسید:

\begin{codebox}{java}
public Student() {}
\end{codebox}

دقت کنید که لازم نیست کانستراکتور قبلی خود را پاک کنید! هر کلاس، می‌تواند
چندین کانستراکتور داشته باشد. شما می‌توانید این دو کانستراکتور را در کنار
هم بنویسید و کدتان همچنان کار می‌کند.

\section{\lr{field} ها}

متغیر های یک کلاس، یا \lr{field} ها، ویژگی‌های آن کلاس را نشان می‌دهند.
شما تا الآن تعداد خوبی \lr{field} برای کلاس \lr{student} تعریف کردید و
با آن‌ها آشنا شدید:

\begin{codebox}{java}
class Student {
    public int age;
    public String name;
    public String studentID;
    public ArrayList<Double> grades;

    // Other things happening in the Student class

}
\end{codebox}

شما حتی می‌توانید در یک کلاس، \lr{field} هایی از جنس همان کلاس تعریف
کنید! مثلاً، می‌توانید فیلد \lr{friend} را برای \lr{\texttt{Student}} ، از
جنس خود \lr{\texttt{Student}} تعریف کنید:

\begin{codebox}{java}
class Student {
    // Other fields
    
    public Student friend;

    // Other things happening in the Student class

}
\end{codebox}

حالا، یک \lr{\texttt{student}} جدید به اسم شهرام توی کدتان تعریف کنید:

\begin{codebox}{java}
var shahram = new Student();
\end{codebox}

برای الآن، بدون این که فیلدهای شهرام را مقداردهی کنید، سعی کنید تا
اطلاعات او را چاپ کنید:

\begin{codebox}{java}
System.out.println("Shahram: ");  
System.out.println("\t+ Name: " + shahram.name);  
System.out.println("\t+ Age: " + shahram.age);  
System.out.println("\t+ StudentID: " + shahram.studentID);
\end{codebox}

کدتان را اجرا کنید، خروجی زیر را برای شهرام می‌بینید:

\begin{codebox}{text}
Shahram:  
    + Name: null  
    + Age: 0  
    + StudentID: null
\end{codebox}

می‌بینید که علی رغم آن که شما به شهرام اسم و سن و شماره دانشجویی ندادید،
خود جاوا یک سری مقدار به آن‌ها اختصاص داده است. مقدار دیفالت جاوا برای
فیلدهایی که مقداردهی نشده‌اند، به این شکل است:

\begin{itemize}
\tightlist
\item
  \textbf{متغیرهای عددی (مثل \lr{\texttt{int}} ، \lr{\texttt{float}} و
  امثال اون‌ها):} مقدار \lr{\texttt{0}} را به خود می‌گیرند.
\item
  \textbf{متغیرهای \lr{\texttt{char}} :} مقدار
  \lr{\texttt{'\textbackslash 0'}} را به خود می‌گیرند.
\item
  \textbf{متغیرهای \lr{\texttt{boolean}} :} مقدار \lr{\texttt{false}} به
  خود می‌گیرند.
\item
  \textbf{متغیرهایی از جنس \lr{reference type} :} مقدار
  \lr{\texttt{null}} به خود می‌گیرند.
\end{itemize}

کلیدواژه‌ی \lr{\texttt{null}} ، یکی از کلیدواژه‌های خاص جاواست که در واقع
یک مقدار نیست، بلکه نشان‌دهنده‌ی آن است که یک متغیر از جنس
\lr{reference type} ، هنوز مقداری به خودش نگرفته است. اگر توی کد
قبلیمان، خط زیر را بنویسیم:

\begin{codebox}{java}
if (shahram.studentID == null) {  
    System.out.println("Shahram does not have a studentID");  
}
\end{codebox}

می‌بینید که پیام زیر چاپ می‌شود:

\begin{codebox}{text}
Shahram does not have a studentID
\end{codebox}

جلوتر، با این کلید‌واژه بهتر آشنا می‌شویم.

\subsection{فیلدهای \lr{static}}

بعضی فیلدها( ویژگی‌ها)، متعلق به هیچ \lr{object} ای نیستند، ولی بی‌ربط به
خود \lr{\texttt{class}} هم نیستند. مثلا در کلاس \lr{\texttt{Student}}
ویژگیِ «تعداد کل دانشجوها» متعلق به هیچ کدام از قلی، ممد یا سلطان نیست،
ولی به کلاس \lr{\texttt{Student}} ربط دارد.

به این ویژگی‌ها، ویژگی‌های \lr{\texttt{static}} گویند. آن‌ها به خود
\lr{class} مرتبط هستند و بین تمام \lr{instance} های آن \lr{class}
مشترکند. فیلد\lr{\texttt{static}} زیر را برای دانشجوها تعریف کنید:

\begin{codebox}{java}
class Student {
    public static int totalNumberOfStudents = 0;

    // other stuff

}
\end{codebox}

سپس در همه‌ی \lr{constructor} هایی که برای \lr{\texttt{Student}}
نوشته‌اید، به مقدار آن یکی اضافه کنید. با این کار، با ساخت هر دانشجو،
تعداد کل دانشجوها یکی زیاد می‌شود:

\begin{codebox}{java}
public Student() {
    totalNumberOfStudents++;  
}

public Student(int newStudentAge, String newStudentName, String newStudentID) {
    age = newStudentAge;
    name = newStudentName;
    studentID = newStudentID;
    grades = new ArrayList<>();
    
    totalNumberOfStudents++;  
}
\end{codebox}

حالا، کد زیر را در \lr{main} بنویسید:

\begin{codebox}{java}
public static void main(String[] args) {
    var gholi = new Student();  
    System.out.println("Current number of students: " + Student.totalNumberOfStudents);

    var mamad = new Student();  
    System.out.println("Current number of students: " + Student.totalNumberOfStudents);

    var javad = new Student();  
    System.out.println("Current number of students: " + Student.totalNumberOfStudents);  
}
\end{codebox}

همچین خروجی‌ای می‌بینید:

\begin{codebox}{text}
Current number of students: 1  
Current number of students: 2  
Current number of students: 3
\end{codebox}

می‌بینید که ما برای دسترسی به \lr{\texttt{totalNumberOfStudents}} ، از
خود کلاس \lr{\texttt{Student}} استفاده کرده‌ایم. می‌توانستید با کد زیر، از
هر کدام از \lr{instance} های \lr{\texttt{student}} نیز به آن دسترسی پیدا
کنید، ولی کار چندان خوبی نیست:

\begin{codebox}{java}
System.out.println("Current number of students: " + gholi.totalNumberOfStudents);
\end{codebox}

\section{\lr{method} ها}

تا اینجای کار، با کلاس‌های نسبتاً ساده‌ای سر و کار داشتید. اما جاهای مختلف
به «رفتار کلاس» یا این ایده که کلاس ما کاری انجام بدهد اشاره کردیم،
اینجاست که متدها وارد عمل می‌شوند: به طور کلی زمانی که بخواهید در کدتان
تصمیمی بگیرید و یا عملیات منطقی یا به طور کلی کاری انجام بدهید، باید از
متدها استفاده کنید. متدها آنقدر مهم هستند که حتی در اولین مواجه‌\-تان با
جاوا از متد \lr{\texttt{main}} استفاده کردید و داخل آن کد خود را نوشتید!
در این بخش قرار است دقیق‌\-تر و کامل‌\-تر با متدها آشنا بشوید. کد زیر، یک
مثال ساده از یک متد است:

\begin{codebox}{java}
public class Refrigerator {
    int numberOfBananas;

    public void getBananas(int n) {
        boolean enoughBananas = numberOfBananas >= n;
        if (enoughBananas) {
            numberOfBananas -= n;  
            System.out.println(
                    "You took " + n + " bananas out of your fridge!"
            );  
        } else {  
            System.out.println(
                    "You don't have that many bananas in your fridge!"
            );  
        }  
    }  
}
\end{codebox}

در این مثال، یک کلاس \lr{\texttt{Refrigerator}} داریم که یک فیلد از نوع
\lr{\texttt{int}} به نام \lr{\texttt{numberOfBananas}} دارد. همچنین یک
متد دارد که نوع خروجی‌اش \lr{\texttt{void}} است (خروجی ندارد) و یک ورودی
( \lr{argument )} از نوع \lr{\texttt{int}} دارد. با استفاده از این متد
می‌\-توانید از توی یخچالتان موز بردارید! حالا خودتان متدی اضافه کنید که
با آن بتوانید توی یخچالتان موز بگذارید. متدتان احتمالاً چیزی شبیه به این
می‌شود:

\begin{codebox}{java}
public void putBananas(int n) {    
    numberOfBananas += n;  
    System.out.println("You put " + n + " bananas in your fridge!");  
}
\end{codebox}

اینجا، متدمان تعداد مشخص و ثابتی ورودی دارد ( یکی )؛ اما می‌\-توانید
متدهایی تعریف کنید که تعداد ورودی\-‌هایشان ثابت نباشد. در این درس نیازی
به یادگیری این کار ندارید اما در صورت علاقه می\-‌توانید کلمه‌ی ``varargs''
را جستجو کنید یا از
\href{https://www.geeksforgeeks.org/variable-arguments-varargs-in-java/}{این
لینک} راجع به آن بخوانید.

\subsection{متغیرهای محلی \lr{(local variables)}}

متد \lr{\texttt{getBananas()}} که در مثال بخش قبل تعریف کردیم، قبل از هر
چیزی چک می‌\-کند که توی یخچال به اندازه‌ی کافی موز وجود داشته باشد و این
را در یک متغیر محلی به اسم \lr{\texttt{enoughBananas}} ذخیره می‌\-کند.
متغیرهای محلی موقتی هستند و فقط توی همان متدی که تعریف شده‌اند قابل
استفاده‌اند. این متغیرها وقتی متد صدا زده می‌شود ساخته می‌شوند و معمولاً بعد
از تمام شدن متد از بین می‌روند. همچنین از بیرون متد نمی‌توانید به آنها
دسترسی داشته باشید. برای آن که خودتان ببینید، در همین کلاس
\lr{\texttt{Refrigerator}} سعی کنید در متد \lr{\texttt{putBananas}} از
متغیر \lr{\texttt{enoughBananas}} استفاده کنید؛ برای مثال:

\begin{codebox}{java}
public void putBananas(int n) {  
    System.out.println(enoughBananas);
    numberOfBananas += n;  
    System.out.println("You put " + n + " bananas in your fridge!");  
}
\end{codebox}

احتمالاً می‌\-\-بینید که \lr{\texttt{enoughBananas}} قرمز شده است. موس را
ببرید روی آن ؛ با همچین چیزی مواجه می‌\-شوید:

\begin{figure}[H]
\centering
\includegraphics[width=0.8\textwidth]{images/oop-1-4.png}
\end{figure}

اینجا \lr{Intellij} دارد به شما می\-‌گوید که نمی‌\-تواند
\lr{\texttt{enoughBananas}} را پیدا کند! دلیلش هم همانطور که گفتیم آن
است که \lr{\texttt{enoughBananas}} داخل متد \lr{\texttt{getBananas}}
تعریف شده و مربوط به همان متد است و داخل \lr{\texttt{putBananas}} همچین
متغیری وجود ندارد! حالا یک متد \lr{\texttt{main}} خالی داخل کلاستان
بنویسید و سعی کنید اجرایش کنید:

\begin{codebox}{java}
public class Refrigerator {
    int numberOfBananas;
    
    public void getBananas(int n) {
        boolean enoughBananas = numberOfBananas >= n;
        if (enoughBananas) {
            numberOfBananas -= n;
            System.out.println(
               "You took " + n + " bananas out of your fridge!"
            );  
            } else {
            System.out.println(
                "You don't have that many bananas in your fridge!"
            );  
        }  
    }

    public void putBananas(int n) {  
        System.out.println(enoughBananas);
        numberOfBananas += n;
        System.out.println("You put " + n + " bananas in your fridge!");  
    }
    
    public static void main(String[] args) {  
    }  
}
\end{codebox}

کدتان کامپایل نمی‌شود. با همچین چیزی مواجه می‌\-شوید:

\begin{figure}[H]
\centering
\includegraphics[width=0.8\textwidth]{images/oop-1-5.png}
\end{figure}

اینجا جاوا نتوانسته کد شما را کامپایل کند و به شما می‌گوید که نمی‌\-تواند
\lr{\texttt{enoughBananas}} را پیدا کند. باید هشدارهای \lr{Intellij} را
جدی می\-‌گرفتید!

ورودی‌های یک متد هم جزو متغیرهای محلی آن متد حساب می‌شوند، با این تفاوت که
مقدار اولیه‌شان موقعی که متد صدا زده می‌شود از طرف کسی که متد را فراخوانی
کرده است، مشخص می‌شود.

\subsubsection{مقداردهی اولیه به متغیرهای محلی}

بر خلاف فیلدهای آبجکت که اگر مقداردهیشان نکنید، جاوا برای آنها مقدار
پیش‌\-فرضی قرار می‌دهد، متغیرهای محلی را باید قبل از استفاده کردن مقداردهی
کنید وگرنه خطای کامپایل می\-‌گیرید:

\begin{codebox}{java}
public class SomeClass {
    // instance variables always get default values if  
    // you don't initialize them
    int foo;

    void myMethod() {
    // local variables do not get default values
        int bar;
        foo += 1; // This is ok, foo has the value 0
        bar += 1; // compile-time error, bar is uninitialized
    }

    public static void main(String[] args) {  
        SomeClass something = new SomeClass();  
        something.myMethod();  
    }  
}
\end{codebox}

سعی کنید کد بالا را اجرا کنید. می‌\-بینید که کدتان حتی کامپایل هم نمی‌شود!
باید اول \lr{\texttt{bar}} را مقداردهی کنید:

\begin{codebox}{java}
bar = 99; // This is ok, we're setting bar's initial value  
bar += 1; // Now this calculation is ok
\end{codebox}

البته دقت کنید که لازم نیست حتماً موقع تعریف کردن یک متغیر آن را مقداردهی
کنید؛ بلکه قبل از آن که از آن استفاده کنید باید مقداردهی شده باشد. موضوع
وقتی پیچیده\-تر می‌شود که مقداردهی را داخل یک شرط انجام بدهید:

\begin{codebox}{java}
void myMethod() {
    int bar;
    if (someCondition) {  
        bar = 42;  
    }  
    bar += 1; // Still a compile-time error, bar may not be initialized
}
\end{codebox}

در این مثال، \lr{\texttt{bar}} فقط در صورتی مقداردهی می‌شود که شرط
\lr{\texttt{someCondition}} برقرار باشد. یعنی همچنان ممکن است که قبل از
خط \lr{\texttt{bar += 1}} ، متغیر \lr{\texttt{bar}} مقداردهی نشده باشد.
کامپایلر به شما اجازه نمی‌دهد همچین کاری بکنید و این کد هم خطای کامپایل
می‌دهد.

برای حل این مشکل، چند راه حل وجود دارد. می\-‌توانید متغیر را قبل از بلوک
شرط مقداردهی کنید، یا استفاده‌ای که از متغیر می‌\-کنید را هم داخل بلوک شرط
ببرید، یا می‌\-توانید با توجه به برنامه\-‌ای که دارید می‌\-نویسید، به نحوی
مطمئن بشید که متغیر قبل از مقداردهی استفاده نمی‌شود. برای مثال، توی کد
بالا می‌\-توانید \lr{\texttt{bar}} را هم در بلوک \lr{if} و هم در بلوک
\lr{\texttt{else}} مقداردهی کنید یا در صورتی که
\lr{\texttt{someCondition}} برقرار نبود، متد را تمام کنید و
\lr{\texttt{return}} کنید:

\begin{codebox}{java}
void myMethod() {
    int bar;
    if (someCondition) {  
        bar = 42;  
    } else {
        return;  
    }  
    bar += 1; // This is ok!
}
\end{codebox}

در این کد، یا \lr{\texttt{bar}} مقداردهی می‌شود و بعد از آن استفاده
می‌شود، یا کلا متد قطع می‌شود و \lr{\texttt{return}} می‌\-کند. جاوا این را
ازتون می\-‌پذیره!

حالا چرا اصلاً جاوا انقدر روی این موضوع حساس است؟ یکی از متداول‌ترین
مشکلاتی که در زبان‌هایی مثل \lr{C} و \lr{C++} به وجود می‌آید آن است که
یادتان می‌رود متغیری را مقداردهی کنید. توی این زبان‌\-ها، متغیرهای
مقداردهی نشده، مقادیر ظاهراً رندومی اختیار می\-‌کنند و این می‌\-تواند
دردسرساز شود و باعث سخت‌تر شدن پروسه‌ی دیباگ برنامه\-‌ها بشود. جاوا با
مجبور کردن شما به مقداردهی به متغیرها، باعث جلوگیری از این مشکل می‌شود.

\subsection{\lr{Shadowing}}

وقتی که نام یک متغیر محلی یا یک ورودی متد با نام یکی از فیلد های کلاسمان
یکی باشد، آن متغیر محلی، اصطلاحا روی آن فیلد ``سایه می‌\-اندازد'' و جلوی
دسترسی ما به آن فیلد را می‌\-گیرد، زیرا در این حالت جاوا نمی‌داند که ما
هنگام صدا زدن آن نام، قصد استفاده از متغیر محلی را داریم یا فیلد متد را،
پس به طور پیش‌\-فرض استفاده از متغیر محلی را اولویت می‌دهد. شاید فکر کنید
این مشکل به ندرت پیش می‌آید، ولی \lr{shadowing} اتفاق نسبتاً متداولی است،
مخصوصا وقتی که متغیرهایمان اسم‌های متداولی داشته باشند. مثال زیر را
ببینید:

\begin{codebox}{java}
public class Car {
    double x;
    double y;
    
    public void moveTo(double x, double y) {  
        System.out.println("The car is moving to " + x + ", " + y);  
    }  
}
\end{codebox}

اینجا ما یک کلاس به نام \lr{\texttt{Car}} داریم که فعلاً فقط دو فیلد برای
مختصات دارد ( \lr{\texttt{x}} و \lr{\texttt{y}} ). یک متد
\lr{\texttt{moveTo}} برای آن تعریف کرده‌ایم که قرار است ماشین را برای ما
حرکت دهد. همانگونه که می‌\-بینید، فعلاً متد \lr{\texttt{moveTo}} فقط دارد
\lr{\texttt{x}} و \lr{\texttt{y}} را چاپ می‌\-کند. اما این
\lr{\texttt{x}} و \lr{\texttt{y}} ، کدام \lr{\texttt{x}} و
\lr{\texttt{y}} هستند؟ اگر مثلاً مختصات ماشین ما \lr{\texttt{(3,4)}} باشد
و ما متد \lr{\texttt{moveTo}} را روی ماشین صدا بزنیم و به آن مقادیر
\lr{\texttt{(6,7)}} را بدهیم، چه چیزی چاپ می‌شود؟ خودتان امتحان کنید! توی
همین کلاس یک متد \lr{\texttt{main}} بنویسید، در آن یک آبجکت جدید از
\lr{\texttt{Car}} بسازید، به آن \lr{\texttt{x}} و \lr{\texttt{y}} بدهید
و متد \lr{\texttt{moveTo}} را رویش صدا بزنید.

همانطور که می‌\-بینید، \lr{\texttt{moveTo}} همان مقادیری را چاپ می\-‌کند
که به آن ورودی دادیم؛ اما اگر بخواهیم ماشین را حرکت بدهیم، باید بتوانیم
مختصاتش را تغییر دهیم، ولی چطور می\-‌توانیم به فیلدهای \lr{\texttt{x}} و
\lr{\texttt{y}} که مربوط به آبجکتمان هستند دسترسی پیدا کنیم؟

\subsubsection{\lr{this}}

هروقت نیاز دارید که صریحاً به آبجکتی که تویش هستیم یا یکی از اعضای آن
اشاره کنید، می‌\-توانید از کلیدواژه‌ی \lr{\texttt{this}} استفاده کنید.
بیاید دوباره با مثال \lr{\texttt{moveTo}} ببینیم:

\begin{codebox}{java}
public class Car {
    double x;
    double y;
    double gas;
    
    public void moveTo(double x, double y) {
        double distance = Math.sqrt(  
            (this.x - x) * (this.x - x) + (this.y - y) * (this.y - y)  
        );
        if (5 * distance > gas) {  
            System.out.println("Not enough gas!");  
        } else {
            this.x = x;
            this.y = y;
            gas -= 5 * distance;  
            System.out.println("The car is moving to " + x + ", " + y);  
        }  
    }  
}
\end{codebox}

اینجا، ابتدا فاصله‌ای که قرار است طی بشود را حساب کرده و توی متغیر محلی
\lr{\texttt{distance}} ریختیم. همانطور که می\-‌بینید، برای دسترسی به
\lr{\texttt{x}} و \lr{\texttt{y}} مربوط به آبجکت ( مختصات فعلی ماشین )،
از \lr{\texttt{this.x}} و \lr{\texttt{this.y}} استفاده کردیم.
\lr{\texttt{this}} در واقع به همان آبجکتی که تویش هستیم اشاره می\-‌کند.

دقت کنید که یک فیلد \lr{\texttt{gas}} نیز به \lr{\texttt{Car}} اضافه
کردیم که قرار است مقدار بنزین ماشین را نشان بدهد. در ادامه‌ی متد، ابتدا
مطمئن می‌\-شویم که ماشین به اندازه‌ی کافی بنزین دارد و سپس ماشین را حرکت
می\-‌دهیم. می‌\-بینید که برای دسترسی به \lr{\texttt{gas}} از
\lr{\texttt{this}} استفاده نکردیم؛ چرا که اشاره به آبجکتی که در آن هستیم
به طور ضمنی برقرار است و \lr{\texttt{gas}} و \lr{\texttt{this.gas}}
اینجا یک چیز هستند.

استفاده از \lr{\texttt{this}} برای دسترسی به فیلدهایی که روی آنها سایه
افتاده، روش مرسومی است و باعث می\-‌شو بتوانیم از اسم\-‌های متداولی که برای
متغیرهای مختلف وجود دارد ( مثلاً \lr{\texttt{x}} و \lr{\texttt{y}} برای
مختصات ) استفاده کنیم و لازم نباشد هر بار دنبال اسمی جدید برای
متغیرهایمان باشیم. علاوه بر این، هر جای دیگری که بخواهید به آبجکتی که
داخلش هستید اشاره کنید، می‌توانید از \lr{\texttt{this}} استفاده کنید.
مثلاً وقتی بخواهید خود آبجکت را به‌عنوان ورودی به یک متد بفرستید.

\subsection{متدهای استاتیک}

متدهای استاتیک \lr{(static methods)} ، مثل فیلدهای استاتیک، به خود کلاس
تعلق دارند، نه به آبجکت‌های مستقلی که ما از روی آن کلاس می‌سازیم. اما این
یعنی چه؟ اول از همه، متدهای استاتیک خارج از آبجکت‌ها وجود دارند و برای
صدا زدنشان لازم نیست آبجکتی وجود داشته باشد؛ شما می‌توانید اسم کلاس را
بنویسید و با عملگر نقطه متدهای استاتیک را صدا بزنید. قبلاً از متدهای
استاتیک زیاد استفاده کرده‌اید، مثلاً برای مرتب کردن آرایه‌ها از
\lr{\texttt{Arrays.sort()}} استفاده می‌کردید در حالی که آبجکتی از کلاس
\lr{\texttt{Arrays}} نساخته بودید و مستقیماً متد \lr{\texttt{sort}} را
روی کلاس \lr{\texttt{Arrays}} صدا می‌زدید؛ این کار قابل انجام است چون
\lr{\texttt{sort}} یک متد استاتیک است.

دوباره از \lr{\texttt{Arrays.sort()}} استفاده کنید یا جایی که آن را
نوشته‌اید را بیاورید، بعد روی آن راست کلیک کنید و گزینه‌ی \lr{Go To} و بعد
\lr{Declaration or Usages} را انتخاب کنید:

\begin{figure}[H]
\centering
\includegraphics[width=0.8\textwidth]{images/oop-1-6.png}
\end{figure}

با همچین کدی مواجه می‌شوید:

\begin{figure}[H]
\centering
\includegraphics[width=0.8\textwidth]{images/oop-1-7.png}
\end{figure}

همان‌طور که می‌بینید، پشت این متد از کلیدواژه‌ی \lr{\texttt{static}}
استفاده شده. این نشان می‌دهد که این متد، یک متد استاتیک است. حال که
می‌دانید متدهای استاتیک چگونه تعریف می‌شوند، بیایید برای کلاس
\lr{\texttt{Car}} که تا الان داشتیم یک متد استاتیک تعریف کنیم:

\begin{codebox}{java}
public class Car {
    public static final int SUV = 0;
    public static final int SEDAN = 1;
    public static final int HATCHBACK = 2;
    
    double x;
    double y;
    double gas;
    int model;

    public static String[] getCarModels() {
        return new String[]{"SUV", "SEDAN", "HATCHBACK"};  
    }
    // ...

\end{codebox}

اینجا، اول از همه یک فیلد جدید به ماشین‌هایمان اضافه کردیم به اسم
\lr{\texttt{model}} که مدل ماشین ما را نشان می‌دهد: ماشین ما می‌تواند
شاسی‌بلند ( \lr{\texttt{model = 0}} ) ، سواری ( \lr{\texttt{model = 1}} )
یا هاچ‌بک ( \lr{\texttt{model = 2}} ) باشد. برای راحتی، مدل‌های مختلف
ماشین را به صورت فیلد‌های \lr{\texttt{static final}} تعریف کردیم. حالا
فرض کنید به اسم این مدل‌ها به صورت \lr{\texttt{String}} نیاز داشته باشیم،
می‌توانیم مثل بالا یک متد استاتیک تعریف کنیم که این اطلاعات را به ما
بدهد. دقت کنید که مدل‌های مختلف ماشین‌ها هیچ ارتباطی به یک ماشین خاص یا یک
آبجکت مشخص از نوع \lr{\texttt{Car}} ندارند و به‌طور کلی برای همه ماشین‌ها
یکسان هستند. به همین خاطر، استفاده از فیلد‌ها و متد‌های استاتیک بهترین
انتخاب است.

اصلی‌ترین کاربرد متدهای استاتیک، برای تعریف متدهای کمکی است؛ متدهایی که
یا مستقل از آبجکت‌ها کار می‌کنند، یا روی آبجکت‌هایی که از آن کلاس (یا حتی
کلاس‌های دیگر) می‌سازیم، عملی انجام می‌دهند و منطقشان به یک \lr{instance}
خاص تعلق ندارد.

حالا بیاید یک متد استاتیک دیگر برای \lr{\texttt{Car}} بنویسیم:

\begin{codebox}{java}
public void printModelsCount(ArrayList<Car> list) {
    int suvCount = 0;
    int sedanCount = 0;
    int hatchbackCount = 0;
    for (Car car : list) {
        switch (car.model) {
            case 0:  
                suvCount++;
                break;
            case 1:  
                sedanCount++;
                break;
            case 2:
                hatchbackCount++;
                break;  
        }  
    }
    System.out.println("SUV: " + suvCount);
    System.out.println("SEDAN: " + sedanCount);
    System.out.println("HATCHBACK: " + hatchbackCount);  
}
\end{codebox}

این متد، یک \lr{ArrayList} از ماشین‌ها می‌گیرد، تعداد مدل‌های مختلف ماشین‌ها
را می‌شمارد و چاپ می‌کند. همان‌طور که می‌بینید کارکرد این متد هیچ ربطی به یک
\lr{instance} خاص از \lr{\texttt{Car}} ندارد و به همین دلیل استاتیک
تعریف می‌شود.

مثال خوب دیگری برای کاربرد متدهای استاتیک، کلاس \lr{\texttt{Math}} است.
این کلاس قرار است مجموعه‌ای از عملیات‌های ریاضی باشد؛ به همین دلیل تمام
متد‌های کلاس \lr{\texttt{Math}} استاتیک هستند. البته \lr{\texttt{Math}}
یک مرحله فراتر می‌رود، شما نمی‌توانید اصلاً آبجکتی از \lr{\texttt{Math}}
بسازید! اصلاً این که یک آبجکت از روی \lr{\texttt{Math}} بسازید، معنی
ندارد و نیازی به آن نیست. این کلاس صرفاً قرار است که مجموعه‌ای از متدها و
متغیرها برای انجام عملیات ریاضی باشد. شما چند ``ریاضی'' مختلف ندارید که
بخواهید instance‌ های مختلفی از \lr{\texttt{Math}} بسازید!

حالا سعی کنید در یکی از متدهای استاتیک \lr{\texttt{Car}} از یکی از
فیلدها یا متدهای غیر استاتیک \lr{\texttt{Car}} استفاده کنید. مثلاً:

\begin{codebox}{java}
public static String[] getCarModels() {  
    System.out.println(model);
    return new String[]{"SUV", "SEDAN", "HATCHBACK"};  
}
\end{codebox}

اگر موستان را ببرید روی \lr{\texttt{model}} یا سعی کنید کلاس
\lr{\texttt{Car}} را کامپایل کنید و جایی از آن استفاده کنید، با همچین
خطاهایی مواجه می‌\-شوید:

\begin{figure}[H]
\centering
\includegraphics[width=0.8\textwidth]{images/oop-1-8.png}
\end{figure}

\begin{figure}[H]
\centering
\includegraphics[width=0.8\textwidth]{images/oop-1-9.png}
\end{figure}

همونطور که می\-بینید، جاوا دارد به شما می‌گوید که نمی‌\-توانید یک فیلد غیر
استاتیک مثل \lr{\texttt{model}} را در یک متد استاتیک استفاده کنید. شما
می‌\-توانید متد \lr{\texttt{getCarModels}} را بدون داشتن هیچ آبجکتی صدا
بزنید؛ پس فیلد \lr{\texttt{model}} ی که سعی دارید از آن استفاده کنید،
مربوط به کدام آبجکت است؟ از آن‌جایی که متدهای استاتیک مربوط به کلاس هستند
و از آبجکت‌\-ها جدا هستند، طبیعی است که نمی‌\-توانند به فیلدها و متدهای
عادی که مربوط به هر آبجکت هستند دسترسی داشته باشند و فقط به متد ها و
متغیر های استاتیک دسترسی دارند.

\subsection{\lr{Method overloading}}

قابلیت \lr{Method overloading} ، این قابلیت است که شما چند متد را با یک
اسم، ولی با نوع و تعداد ورودی متفاوت در یک کلاس تعریف کنید؛ وقتی که متد
را صدا می‌زنید، کامپایلر با توجه به نوع ورودی، متد درست را انتخاب می‌کند و
اجرا می‌کند.

قابلیت \lr{Method overloading} ، بسیار قدرتمند و پرکاربرد است. ایده اصلی
این است که متدهایی درست کنیم که روی ورودی‌های مختلف، کارهای یکسانی انجام
بدهند. با این کار می‌توانید این توهم را ایجاد کنید که یک متد می‌تواند روی
انواع مختلفی از ورودی‌ها کار کند. متد \lr{\texttt{println()}} که از اولین
جلسه با آن کار کردید، مثال خیلی خوبی از \lr{method overloading} است؛ شما
می‌توانید به \lr{\texttt{println()}} هر ورودی دلخواهی بدهید و آن به نحوی
یک نمایش متنی از آن ورودی را برای شما چاپ می‌کند. در زبان‌هایی که
\lr{method overloading} ندارند، کار سخت‌تر می‌شود. مثلاً برای چاپ چیزهای
مختلف باید متدهای مختلف با اسم‌های مختلف تعریف کنیم و در آن صورت، این
مسئولیت روی دوش شما می‌افتد که متد درست را انتخاب کنید. بیایید یک مثال
دیگر از \lr{method overloading} ببینیم:

\begin{codebox}{java}
public class Sum {
    // Overloaded sum(). This sum takes two int parameters
    public int sum(int x, int y) {
        return (x + y);  
    }

    // Overloaded sum(). This sum takes three int parameters
    public int sum(int x, int y, int z) {
        return (x + y + z);  
    }

    // Overloaded sum(). This sum takes two double  
    // parameters
    public double sum(double x, double y) {
        return (x + y);  
    }

    public static void main(String[] args) {  
        Sum s = new Sum();  
        System.out.println(s.sum(10, 20));  
        System.out.println(s.sum(10, 20, 30));  
        System.out.println(s.sum(10.5, 20.5));  
    }  
}
\end{codebox}

همان‌طور که می‌بینید، اینجا سه تا متد با نام یکسان \lr{\texttt{sum}}
داریم، ولی ورودی‌های آن‌ها متفاوت است. هر سه‌تای این متدها دارند عمل جمع
کردن را انجام می‌دهند، ولی یکی دو \lr{\texttt{double}} را، یکی دو‌
\lr{\texttt{int}} را و یکی سه \lr{\texttt{int}} را جمع می‌کند! به غیر از
نوع ورودی\-ها و تعدادشون، با تغییر دادن ترتیب ورودی\-ها هم میشود متدها
را \lr{overload} کرد:

\begin{codebox}{java}
class Student {
    // Method 1
    public void getStudentInfo() {  
        System.out.println("Name :" + name + " "
                + "ID :" + roll_ studentID);  
    }

    // Method 2
    public void getStudentInfo (String name) {
        // Again printing name and id of person
        System.out.println("ID :" + studentID + " "
                + "Name :" + name);  
    }  
}
\end{codebox}

بعد از این که با مباحث مربوط به \lr{polymorphism} و متدهای \lr{override}
شده آشنا شدید، دوباره به \lr{method overloading} برخواهیم گشت.

\section{\lr{Reference type} ها}

همان‌طور که تا الآن به خوبی می‌دانید، در جاوا، \lr{type} ها به دو دسته‌ی
\lr{primitive type} و \lr{reference type} تقسیم‌بندی می‌شوند.
\lr{primitive type} ها، تایپ‌های بسیار ساده‌ای مثل \lr{\texttt{int}} ،
\lr{\texttt{char}} ، \lr{\texttt{boolean}} و امثال آن‌ها هستند. فهرست
کامل آن‌ها در لیست زیر آمده است:

\begin{table}[H]
\centering
\setLTR
\begin{tabular}{lll}
\hline
\textbf{\lr{Type}} & \textbf{\lr{Definition}} & \textbf{\lr{Approximate range or precision}} \\
\hline
\texttt{boolean} & \lr{Logical Value} & \texttt{true} \lr{or} \texttt{false} \\
\texttt{char} & \(16\)\lr{-bit, Unicode character} & \(64k\) \lr{characters} \\
\texttt{byte} & \(8\)\lr{-bit, signed integer} & \(-128\) \lr{to} \(127\) \\
\texttt{short} & \(16\)\lr{-bit, signed integer} & \(-32, 768\) \lr{to} \(32,676\) \\
\texttt{int} & \(32\)\lr{-bit, signed integer} & \(-2.1\mathrm{e}9\) \lr{to} \(2.1\mathrm{e}9\) \\
\texttt{long} & \(64\)\lr{-bit, signed integer} & \(-9.2\mathrm{e}18\) \lr{to} \(9.2\mathrm{e}18\) \\
\texttt{float} & \(32\)\lr{-bit, IEEE 754, floating-point value} & \(6-7\) \lr{significant decimal places} \\
\texttt{double} & \(64\)\lr{-bit, IEEE} 754 & \(15\) \lr{significant decimal places} \\
\hline
\end{tabular}
\end{table}

هر تایپ دیگری در جاوا، \lr{reference type} است . \lr{\texttt{String}} ،
\lr{\texttt{JFrame}} ، \lr{\texttt{ArrayList}} و حتی تایپ‌هایی مثل
\lr{\texttt{Car}} و \lr{\texttt{Student}} که تا این‌جای کار تعریف کردیم،
همگی \lr{reference type} اند. هر \lr{reference type} ای با یک کلاس تعریف
شده است.

\subsection{تفاوت \lr{reference type} ها و \lr{primitive type} ها}

همان‌طور که می‌دانید، تمام متغیرهای برنامه‌های شما، در حافظه‌ی خاصی به اسم
\lr{RAM} ذخیره می‌شوند. سیستم عامل، متغیرهای شما را در دو بخش متفاوت از
این حافظه، به اسم \lr{stack} و \lr{heap} نگه می‌دارد. با این دو در درس‌های
ساختمان داده و سیستم عامل بیشتر آشنا می‌شوید، ولی فعلأ بدانید که حافظه‌ی
\lr{stack} ، از \lr{heap} سریع‌تر است، ولی در مقابل کمی کم‌حجم‌تر است.
\footnote{البته، در واقعیت ممکن است اینگونه نباشد. \lr{Stack} و
  \lr{Heap} به خودی خود از دیگری سریع‌تر یا حجیم‌تر نیستند و ما در این‌جا
  خیلی ساده‌سازی می‌کنیم. در درس‌های بعدی خود بهتر می‌فهمید که تفاوت این دو
  با هم چیست. برای الآن، فرض کنید چیزی که گفتیم کاملاً درست است.}

متغیرهایی که از جنس \lr{primitive type} تعریف می‌کنید، حجم کمی دارند و
بین \lr{\texttt{1}} تا \lr{\texttt{8}} بایت از حافظه را اشغال می‌کنند. به
همین خاطر، جاوا آن‌ها را در \lr{stack} نگه می‌دارد تا از سرعت بهتر
\lr{stack} استفاده کند و همزمان، \lr{stack} سریع پر نشود. از طرفی، ممکن
است object‌ هایی که در برنامه‌ي خود تعریف می‌کنید -و همیشه از جنس
\lr{reference type} هستند-، حجم بسیار بیشتری داشته باشند. جاوا، اطلاعات
این object‌ها را در \lr{heap} ذخیره می‌کند و در \lr{stack} ، صرفاً یک
اشاره‌گر \lr{(pointer)} به آن‌ها نگه می‌دارد.

مثلاً، برنامه‌ی زیر را در نظر بگیرید:

\begin{codebox}{java}
public class Main {
    public static void main(String[] args) {
        int radius = 5;
        double pi = 3.14;  
        Student changiz = new Student(50, "Changiz", "112");  
    }  
}
\end{codebox}

اگر مموری را حین اجرای این برنامه ببینیم، همچین شکلی دارد:

\begin{figure}[H]
\centering
\includegraphics[width=0.8\textwidth]{images/oop-1-10.png}
\end{figure}

اگر مقدار خود متغیر چنگیز را چاپ کنید:

\begin{codebox}{java}
System.out.println(changiz);
\end{codebox}

همچین خروجی‌ای می‌بینید:

\begin{codebox}{text}
Student@6acbcfc0
\end{codebox}

متغیر \lr{\texttt{changiz}} ، در واقع صرفا یک \lr{pointer} یا
\lr{reference} به یک آبجکت از جنس \lr{Student} است و نه چیزی بیشتر. با
استفاده از اپراتور نقطه \lr{\texttt{(.)}} ، می‌توانید به \lr{field} ها و
\lr{method} های این آبجکت دسترسی داشته باشید.

این اتفاق، یک ساید افکت جالب روی کدهای شما دارد. فرض کنید، یک دانشجوی
دیگر به اسم بنگیز درست کردیم و ان را مساوی با چنگیز قرار دادیم:

\begin{codebox}{java}
Student bangiz = changiz;
\end{codebox}

حالا، بنگیز و چنگیز هر دو به یک نقطه از \lr{heap} اشاره می‌کنند:

\begin{figure}[H]
\centering
\includegraphics[width=0.8\textwidth]{images/oop-1-11.png}
\end{figure}

حالا اگر شما، شماره دانشجویی چنگیز را عوض کنید:

\begin{codebox}{java}
changiz.studentID = "40113";
\end{codebox}

و بعد، شماره دانشجویی چنگیز و بنگیز را چاپ کنید:

\begin{codebox}{java}
System.out.println("Changiz studentID: " + changiz.studentID);  
System.out.println("Bangiz studentID: " + bangiz.studentID);
\end{codebox}

می‌بینید که شماره دانشجویی بنگیز هم عوض شده است!

\begin{codebox}{java}
Changiz studentID: 40113  
Bangiz studentID: 40113
\end{codebox}

عجیب نیست؟ اگر به مموری نگاه کنید، بعد از عوض شدن شماره دانشجویی چنگیز،
همچین وضعیتی دارد:

\begin{figure}[H]
\centering
\includegraphics[width=0.8\textwidth]{images/oop-1-12.png}
\end{figure}

چنگیز، شماره دانشجویی آبجکتی که به ان اشاره می‌کرده را تغییر داده است.
دست بر قضا، بنگیز هم به همین آبجکت اشاره می‌کرده و در نتیجه، شماره
دانشجویی بنگیز هم واقعا عوض شده است.

یک جای دیگر هم اتفاق مشابه‌ای می‌افتد. برای دیدن آن، متد زیر را به کدتان
اضافه کنید:

\begin{codebox}{java}
public static void resetID(Student student) {  
    student.studentID = "00000000";  
}
\end{codebox}

این متد، یک \lr{\texttt{Student}} را ورودی می‌گیرد، و شماره دانشجویی‌ او
را دستکاری می‌کند. حالا با استفاده از کد زیر، چنگیز را به این متد ورودی
بدهید و بعد، شماره دانشجویی چنگیز و بنگیز را چاپ کنید:

\begin{codebox}{java}
resetID(changiz);  
  
System.out.println("Changiz studentID: " + changiz.studentID);  
System.out.println("Bangiz studentID: " + bangiz.studentID);
\end{codebox}

خروجی، به این شکل است:

\begin{codebox}{text}
Changiz studentID: 00000000  
Bangiz studentID: 00000000
\end{codebox}

شاید بتوانید حدس بزنید که این‌جا چه اتفاقی افتاد. با ورودی دادن چنگیز به
\lr{\texttt{resetID}} ، در واقع شما اشاره‌گرتان به \lr{heap} را به این
تابع ورودی دادید. پس متغیر \lr{\texttt{student}} در
\lr{\texttt{resetID}} و \lr{\texttt{changiz}} و \lr{\texttt{bangiz}} ،
هر سه به یک نقطه از \lr{heap} اشاره می‌کنند و مثل قبل، با تغییر فیلدهای
یکی از آن‌ها، هر سه تغییر می‌کنند. به این نوع ورودی دادن به توابع، اصطلاحاً
\lr{passing by reference} می‌گویند.

مشابه هیچ کدام از این اتفاق‌ها، برای \lr{primitive} type‌ ها نمی‌افتد. چون
همیشه در \lr{stack} نگه‌داری می‌شوند و پوینتری به \lr{heap} برای آن‌ها
نگه‌داری نمی‌شود.

\subsection{کلاس‌های \lr{wrapper} برای \lr{primitive} type‌ها}

هر کدام از \lr{primitive} type‌ ها، تایپی مشابه از جنس
\lr{reference type} هم دارند. این تایپ‌ها، در جدول زیر آمده‌اند:

\begin{table}[H]
\centering
\setLTR
\begin{tabular}{ll}
\hline
\textbf{\lr{Primitive}} & \textbf{\lr{Wrapper}} \\
\hline
\texttt{void} & \texttt{java.long.Void} \\
\texttt{boolean} & \texttt{java.lang.Boolean} \\
\texttt{char} & \texttt{java.lang.Character} \\
\texttt{byte} & \texttt{java.lang.Byte} \\
\texttt{short} & \texttt{java.lang.Short} \\
\texttt{int} & \texttt{java.lang.Integer} \\
\texttt{long} & \texttt{java.lang.Long} \\
\texttt{float} & \texttt{java.lang.Float} \\
\texttt{double} & \texttt{java.lang.Double} \\
\hline
\end{tabular}
\end{table}

نیازی نیست که خیلی درگیر این تایپ‌های جدید باشید، ولی بدانید که وجود
دارند. وقتی که یک آرایه از \lr{int} ها تعریف می‌کنید، از آن‌ها استفاده
می‌کنید:

\begin{codebox}{java}
var arr = new ArrayList<Integer>();
\end{codebox}

اگر دقت کنید، به جای این که بین دو براکت از \lr{\texttt{int}} استفاده
کنید، از \lr{\texttt{Integer}} استفاده کردید. اگر روی آن کلیک راست کنید
و از \lr{Go To} به \lr{Declaration and Usages} بروید، می‌توانید ببینید که
پشت آن یک کلاس است:

\begin{figure}[H]
\centering
\includegraphics[width=0.8\textwidth]{images/oop-1-13.png}
\end{figure}

این کلاس، معادل \lr{reference type} گونه‌ای برای \lr{\texttt{int}} است، و
وقتی با کلاس‌ها سر و کار دارید از آن استفاده می‌شود. اگر خواستید بیشتر
راجع به آن بدانید، به
\href{https://docs.oracle.com/javase/tutorial/java/data/autoboxing.html}{این
داکیومنت} رجوع کنید.

\section{\lr{Garbage Collection}}

قبل از اینکه برویم سراغ \lr{Garbage Collection} ، بیایید اول مفهوم نشت
حافظه \lr{(Memory Leak)} را بررسی کنیم.

\subsection{آشنایی با \lr{Memory Leak}}

در بعضی از زبان‌های برنامه‌نویسی (مثل \lr{C} و \lr{C++)} مسئولیت «آزاد
کردن» حافظه بر عهده خود برنامه‌نویس است. این یعنی شما باید هر وقت که دیگر
به یک شی نیاز نداشتید، خودتان آن حافظه‌ی را آزاد کنید. مثلاً اگر با
استفاده از تابع \lr{\texttt{malloc}} مقداری از حافظه را تخصیص
\lr{(allocate)} دادید، وقتی که دیگر به این حافظه نیاز نداشتید، باید
خودتان با استفاده از تابع \lr{\texttt{free}} آن حافظه را آزاد کنید.
بیایید یک مثال ببینیم:

\begin{codebox}{c}
#include <stdlib.h>

int main() {
    // Allocate memory dynamically
    int *ptr = (int *)malloc(sizeof(int) * 5);  // Allocating memory for 5 integers  
  
    // Use the allocated memory
    for (int i = 0; i < 5; i++) {  
        ptr[i] = i + 1;  
    }
    /*  
     * Doing some stuff with these numbers  
     */  
  
    // Forgetting to free the allocated memory causes a Memory Leak  
    // free(ptr); // If we uncomment this line, the Memory Leak will be avoided.

    return 0;  
}
\end{codebox}

در این برنامه که به زبان \lr{C (!)} نوشته شده، اول به اندازه
\lr{\texttt{5}} متغیر \lr{\texttt{int}} حافظه اشغال می‌کنیم. اشاره‌گر
\lr{\texttt{ptr}} به اولین خانه از این \lr{\texttt{20}} بایت \footnote{با
  این فرض که هر متغیر از جنس \lr{\texttt{int}} ، حافظه‌ای به اندازه
  \lr{\texttt{4}} بایت را اشغال کند.} اشاره می‌کند. حال می‌آییم
\lr{\texttt{5}} تا عدد صحیح را در حافظه ذخیره می‌کنیم. فرض کنید با این
اعداد کارهایی انجام دادیم و الان کارمان با آن‌ها تمام شده است اما فراموش
کردیم که این \lr{\texttt{20}} بایت حافظه را آزاد کنیم! در حقیقت باید با
صدا زدن تابع \lr{\texttt{free}} اعلام می‌کردیم که ما دیگر به این
\lr{\texttt{20}} بایت نیازی نداریم تا این حافظه آزاد بشود.

اگر آزادش نکنید چه می‌شود؟ در این صورت آن حافظه همچنان در اختیار برنامه
قرار دارد و به اصطلاح \lr{Memory Leak} رخ می‌دهد. این یعنی حافظه‌ای که
دیگر به کار نمی‌آید، همچنان در اختیار برنامه باقی می‌ماند و هیچ‌وقت آزاد
نمی‌شود. این موضوع می‌تواند باعث شود که برنامه به مرور زمان حافظه‌ی زیادی
مصرف کند و سیستم دچار مشکلاتی مثل کندی عملکرد یا حتی \lr{crash} بشود.

\subsection{\lr{Garbage Collection:} یک راه‌حل خوب}

زمانی که داریم راجع به \lr{Garbage Collection} در جاوا صحبت می‌کنیم، به
زبان ساده یعنی جاوا خودش می‌رود و حافظه‌ای که دیگر به هیچ‌کار نمی‌آید را
آزاد می‌کند. مثلاً وقتی که شما یک شی را در برنامه می‌سازید و دیگر به آن
نیاز ندارید، جاوا خود به خود این شی را پاک می‌کند. شما اصلاً نیازی نیست که
خودتان حافظه را آزاد کنید، همه چیز به صورت خودکار اتفاق می‌افتد!

شاید بپرسید که چطور این کار انجام می‌شود؟ خب، جاوا از تعدادی الگوریتم
برای این کار استفاده می‌کند. ولی ما آن‌ها را اینجا بررسی نمی‌کنیم. در این
\href{https://youtu.be/Mlbyft_MFYM?si=jhAdnlq12houCo1F}{ویدیوی یوتیوب} و
این
\href{https://docs.oracle.com/en/java/javase/21/gctuning/introduction-garbage-collection-tuning.html}{داک
اوراکل} می‌توانید مطالب بیشتری در رابطه با این موضوع ببینید.

\begin{figure}
\centering
\pandocbounded{\includegraphics[width=0.45\textwidth]{images/oop-1-14.png}}
\caption{\textbf{\lr{Thumbnail:} \emph{Java's
\lr{Garbage Collection Explained - How It Saves your Lazy Programmer Bu}}}}
\end{figure}

به طور کلی این ویژگی جاوا باعث می‌شود که شما تمرکز بیشتری روی منطق
برنامه‌نویسی داشته باشید و دیگر نگران مدیریت دستی حافظه نباشید.

حالا در کد زیر، می‌توانید رفتار \lr{Garbage Collector} را در ارتباط با
آبجکت \lr{\texttt{cuteCat}} بگویید؟

\begin{codebox}{java}
Cat cuteCat = new Cat("Cat 1");  
cuteCat = new Cat("Cat 2");
\end{codebox}

شکل زیر میتواند نمایش خوبی از اتفاقات باشد. در اینجا
\lr{\texttt{cuteCat}} به جایی از حافظه اشاره \footnote{دقت کنید که لفظ
  «اشاره کردن» در جاوا خیلی درست نیست. به خاطر این‌که ما در جاوا
  \lr{\texttt{pointer}} نداریم و مدیریت حافظه را خود جاوا برایمان انجام
  می‌دهد. در جاوا حتی دسترسی مستقیم به حافظه هم نداریم (برخلاف
  \lr{\texttt{C}} ). اما برای مشابهت با زبان \lr{\texttt{C}} ، در اینجا
  هم از لفظ اشاره کردن استفاده کردیم. در جاوا به متغیر هایی مثل
  \lr{\texttt{cuteCat}} ، می گویند \lr{reference variable} .} می‌کند که
در آن یک آبجکت از جنس \lr{\texttt{Cat}} (که فیلد اسمش \lr{\texttt{Cat1}}
است) ذخیره شده. در خط بعدی \lr{\texttt{cuteCat}} به جایی از حافظه اشاره
می‌کند که در آن یک آبجکت از جنس \lr{\texttt{Cat}} (که فیلد اسمش \lr{Cat2}
است) ذخیره شده باشد. اما الان هیچ فلشی به \lr{\texttt{Cat1}} وارد
نمی‌شود، پس \lr{Garbage Collector} آن را آزاد می‌کند.

\begin{figure}[H]
\centering
\includegraphics[width=0.8\textwidth]{images/oop-1-15.png}
\end{figure}

\subsection{\lr{Garbage Collector} حواسش هست!}

\lr{Garbage Collector} آبجکت‌هایی که هنوز نیازشان داریم (بهشان رفرنس
داریم) را پاک نمی‌کند. کد زیر را ببینیم:

\begin{codebox}{java}
class Duck {  
    String name;  
}

public class Main {
    public static void main(String[] args) {  
        Duck duck = createDuck(); // a duck object will be created
        System.out.println(duck.name);  
    }


    public static Duck createDuck() {  
        Duck localDuck = new Duck();  
        localDuck.name = "A White Duck";
        return localDuck;  
    }  
}
\end{codebox}

در این کد ابتدا یک کلاس خیلی ساده به نام \lr{\texttt{Duck}} تعریف کردیم.
متدی به نام \lr{\texttt{createDuck}} نیز تعریف کردیم که ابتدا یک آبجکت
از جنس \lr{\texttt{Duck}} می‌سازد، سپس متغیر \lr{\texttt{name}} را در این
آبجکت مقداردهی می‌کند و در نهایت این آبجکت را به عنوان خروجی متد
برمی‌گرداند. بعد در متد \lr{\texttt{main}} ، سعی می‌کنیم به متغیر
\lr{\texttt{name}} در این آبجکت دسترسی داشته باشیم.

اما نکته کجاست؟ احتمالاً در درس‌های قبلیتان خواندید که \emph{«متغیر‌هایی که
در یک تابع تعریف می‌شوند عمرشان به اندازه‌ی اجرای همان تابع است و بعد از
اتمام اجرای تابع، آن متغیر هم از بین می‌رود.»} پس شاید انتظار داشته باشیم
\lr{Garbage Collector} آبجکت \lr{\texttt{Duck}} را از بین ببرد! اما
واقعیت آن است که \lr{Garbage Collector} \emph{حواسش است} که ما به کدام
آبجکت‌ها هنوز نیاز داریم و نباید پاکشان کند. اینجا هم چون
\lr{\texttt{Duck}} را به عنوان خروجی برگرداندیم، یعنی آن را نیاز داریم،
پس پاکش نمی‌کند.

خروجی کد بالا به صورت زیر است:

\begin{codebox}{text}
A White Duck
\end{codebox}

\subsection{یک نکته در مورد کلاس‌ها}

در هر فایل جاوا (فایل با پسوند \lr{\texttt{.java}})، می‌توانیم حداکثر یک
کلاس \lr{\texttt{public}} داشته باشیم، همچنین اسم این کلاس
\lr{\texttt{public}} باید حتما با اسم فایل یکی باشد. مثال درست زیر را
ببینید:

\begin{codebox}{java}
// MyClass.java
public class MyClass {
    public void sayHello() {  
        System.out.println("Hello, Java!");  
    }  
}
\end{codebox}

کد‌های زیر نادرست هستند:

\begin{codebox}{java}
// MyFile.java
public class MyClass { } // Error! The class name does not match the file name.
\end{codebox}

\begin{codebox}{java}
// MyFile.java
public class MyClass { }

public class AnotherClass { } // Error! Only one public class is allowed.
\end{codebox}

\section{\lr{Packages}}

\subsection{نیاز به منظم کردن فایل‌ها}

در برنامه‌هایی که در جاوا می‌نویسیم، همیشه از کلاس ها یا اینترفیس ها
\footnote{با اینترفیس‌ها بعداً آشنا می‌شوید. فعلاً کاری به آن‌ها نداریم.}
استفاده می کنیم. برنامه‌ی ساده زیر که در کلاس \lr{Sample} نوشته شده را
ببینید:

\begin{codebox}{java}
public class Sample {
    public static void main(String[] args) {  
        System.out.println("Hello World!");  
    }  
}
\end{codebox}

ولی برنامه‌های پیچیده‌تر ممکن است از صد‌ها کلاس تشکیل شده باشند. اگر همه‌ی
این کلاس‌ها را بدون هیچ نظمی کنار هم قرار بدهیم باعث می‌شود برنامه‌مان
ناخوانا باشد و خودمان هم گیج می‌شویم.

احتمالاً یکی از اولین چیز‌هایی که برای منظم کردن فایل‌ها به ذهن ما می‌رسد
استفاده از پوشه‌هاست. و این دقیقاً همان امکانی است که جاوا برای منظم کردن
کلاس‌های برنامه‌مان برای ما فراهم کرده: ایجاد \lr{\texttt{package}}‌های
مختلف.

\subsection{پکیج چیست؟}

پکیج‌ها در جاوا مثل پوشه‌هایی هستند که کدهای برنامه‌نویسی (شامل کلاس‌ها و
اینترفیس‌ها) را در آن‌ها قرار می‌دهیم تا همه چیز منظم و مرتب باشد. وقتی که
برنامه‌های پیچیده‌تر را می‌نویسیم، تعداد کلاس‌ها زیاد می‌شود و اینجاست که
پکیج‌ها به کمکمان می‌آیند تا بتوانیم این کلاس‌ها را دسته‌بندی کنیم.

\begin{figure}[H]
\centering
\includegraphics[width=0.8\textwidth]{images/oop-1-16.png}
\end{figure}

\subsection{چطور یک پکیج ایجاد کنیم؟}

اول یک پروژه به اسم \lr{\texttt{PackageDemo}} ایجاد کنید. حال روی پوشه‌ی
\lr{\texttt{src}} راست کلیک کنید و از نوار \lr{New} ، گزینه‌ی
\lr{Package} را انتخاب کنید.

\begin{figure}[H]
\centering
\includegraphics[width=0.8\textwidth]{images/oop-1-17.png}
\end{figure}

حالا اسم پکیج را وارد کنید. مثلا در اینجا \lr{\texttt{edu.ap.animals}} :

\begin{figure}[H]
\centering
\includegraphics[width=0.8\textwidth]{images/oop-1-18.png}
\end{figure}

به همین شکل پکیج \lr{\texttt{edu.ap.vehicles}} را هم بسازید.

الآن باید چیزی شبیه به تصویر زیر را در پوشه‌ی \lr{\texttt{src}} داشته
باشید:

\begin{figure}[H]
\centering
\includegraphics[width=0.8\textwidth]{images/oop-1-19.png}
\end{figure}

حال می‌خواهیم چند کلاس به هر کدام از این پکیج‌ها اضافه کنیم. مثلأ کلاس‌های
\lr{\texttt{Dog}} و \lr{\texttt{Cat}} را به پکیج \lr{\texttt{animals}} و
کلاس‌های \lr{\texttt{Bike}} و \lr{\texttt{Car}} را به پکیج
\lr{\texttt{vehicles}} اضافه کنیم. در نهایت باید چیزی شبیه به تصویر زیر
را داشته باشید:

\begin{figure}[H]
\centering
\includegraphics[width=0.8\textwidth]{images/oop-1-20.png}
\end{figure}

پکیج‌هایمان را ساختیم! حالا ببینیم واقعا چه فایل‌هایی ایجاد شده است. اگر
از توی \lr{explorer} به محل ایجاد پروژه‌تان بروید و به ترتیب وارد پوشه‌های
\lr{\texttt{src}} ، \lr{\texttt{edu}} ، \lr{\texttt{ap}} و
\lr{\texttt{animals}} بشوید چیزی شبیه به این‌ها می‌بینید:

\begin{figure}[H]
\centering
\includegraphics[width=0.8\textwidth]{images/oop-1-21.png}
\end{figure}

\begin{figure}[H]
\centering
\includegraphics[width=0.8\textwidth]{images/oop-1-22.png}
\end{figure}

\begin{figure}[H]
\centering
\includegraphics[width=0.8\textwidth]{images/oop-1-23.png}
\end{figure}

\begin{figure}[H]
\centering
\includegraphics[width=0.8\textwidth]{images/oop-1-24.png}
\end{figure}

\begin{figure}[H]
\centering
\includegraphics[width=0.8\textwidth]{images/oop-1-25.png}
\end{figure}

پس دیدیم که وقتی ما داریم یک پکیج تعریف می‌کنیم، واقعا پوشه ایجاد می‌شود!

نهایتأ پکیج \lr{\texttt{edu.ds}} شامل کلاس \lr{\texttt{Stack}} را درست
کنید. در نهایت باید چیزی شبیه به این ببینید:

\begin{figure}[H]
\centering
\includegraphics[width=0.8\textwidth]{images/oop-1-26.png}
\end{figure}

الآن برویم کد \lr{Cat.java} را ببینیم:

\begin{codebox}{java}
package edu.ap.animals;

public class Cat {  
}
\end{codebox}

همان‌طور که می‌بینید، عبارت \lr{\texttt{package edu.ap.animals;}} به
ابتدای این کد اضافه شده است. این خط را نباید پاک کنید، زیرا در آن صورت
جاوا متوجه نمی‌شود \lr{\texttt{Cat}} متعلق به کدام پکیج است و خطای
کامپایل دریافت خواهید کرد.

\subsection{نام‌گذاری متداول پکیج‌ها \lr{(Naming Conventions)}}

\begin{enumerate}
\def\labelenumi{\arabic{enumi}.}
\tightlist
\item
  پکیج‌ها همواره با حروف کوچک نام‌گذاری می‌شوند. مثلاً
  \lr{\texttt{java.util}}
\item
  از کلمات رزرو شده‌ی جاوا استفاده نکنید. کلماتی مثل \lr{\texttt{class}}
  ، \lr{\texttt{public}} ، \lr{\texttt{static}} و\ldots{}
\item
  از زیرخط یا \lr{underscore} \lr{\texttt{\_}} ، خط‌تیره یا \lr{dash}
  \lr{\texttt{-}} ، فاصله یا ~space و کاراکترهای خاص (مثل
  \lr{\texttt{@}} ، \lr{\texttt{\$}} ، \lr{\texttt{\&}}) استفاده نکنید.
  در نام‌گذاری پکیج‌ها فقط مجازیم از \lr{dot (} نقطه \lr{\texttt{.}} )
  استفاده کنیم.

  \begin{enumerate}
  \def\labelenumii{\arabic{enumii}.}
  \tightlist
  \item
    نادرست: \lr{\texttt{com.github.my-awesome-project}}
  \item
    درست: \lr{\texttt{com.github.myawesomeproject}}
  \end{enumerate}
\item
  یک قاعده‌ی دیگر در نام‌گذاری پکیج‌ها، \lr{reverse domain name} است (برعکس
  نوشتن نام دامنه). اگر یک پروژه‌ای متعلق به شرکت یا سازمانی است که دامنه
  ( \lr{domain )} خودش را دارد، آن دامنه را به شکل برعکس می‌نویسیم. مثلاً
  شرکت \lr{Mozilla} (که نام دامنه‌اش \lr{\texttt{mozilla.org}} است) چند
  پکیج در جاوا دارد. یکی از این پکیج‌ها اسمش
  \lr{\texttt{org.mozilla.javascript}} است (البته این قاعده در مورد
  پکیج‌های استاندارد خود جاوا صدق نمی‌کند).
\end{enumerate}

\subsection{کلیدواژه‌ی \lr{Import}}

\subsubsection{یک مثال عملی}

بیایید در کلاس \lr{\texttt{Bike}} یک آبجکت از کلاس \lr{\texttt{Cat}}
بسازیم (دقت کنید که این دو کلاس متعلق به دو پکیج متفاوت هستند):

\begin{codebox}{java}
package edu.ap.vehicles;

public class Bike {
    public static void main(String[] args) {  
        Cat cat = new Cat();  
    }  
}
\end{codebox}

اگر سعی کنیم این کد را اجرا کنیم، موقع کامپایل به مشکل می‌خوریم:

\begin{codebox}{bash}
java: cannot find symbol  
symbol:   class Cat  
  location: class edu.ap.vehicles.Bike
\end{codebox}

مشکل چیست ؟ جاوا نمیتواند کلاسی به اسم \lr{\texttt{Cat}} را پیدا کند!
دلیلش هم این است که جاوا فقط کلاس‌هایی را می‌بیند که توی همین پکیج هستند.

پس باید به نحوی کلاس \lr{\texttt{Cat}} را به کدمان اضافه کنیم. کلیدواژه‌ی
\lr{\texttt{import}} دقیقا برای همین کار است. با استفاده از این
کلیدواژه، ما به جاوا اعلام می‌کنیم که می‌خواهیم این کلاس را به کدمان اضافه
کنیم:

\begin{codebox}{java}
package edu.ap.vehicles;
import edu.ap.animals.Cat;

public class Bike {
    public static void main(String[] args) {  
        Cat cat = new Cat();  
    }  
}
\end{codebox}

اکنون دیگر کدمان کار می‌کند و می‌توانیم یک آبجکت از کلاس \lr{\texttt{Cat}}
بسازیم.

\subsection{در اهمیت پکیج‌ها}

یکی از بزرگترین مزایای جاوا داشتن کتابخانه (library)‌ های بسیار متنوع و
کاربردی است.

می‌خواهید یک \lr{PDF} درست کنید؟ کتابخانه مربوط به آن را
\lr{\texttt{import}} کنید. می‌خواهید با دیتابیس کار کنید؟ کتابخانه مربوط
به آن را \lr{\texttt{import}} کنید \ldots{}

مثلاً یکی از پکیج‌های مهم جاوا، \lr{\texttt{java.lang}} است. این پکیج شامل
کلاس‌های پایه‌ای مثل \lr{\texttt{System}} ، \lr{\texttt{Integer}} ،
\lr{\texttt{Math}} و \lr{\texttt{String}} است.

در شکل زیر می‌توانید یک نمای کلی از پکیج‌های جاوا ببینید:

\begin{figure}[H]
\centering
\includegraphics[width=0.8\textwidth]{images/oop-1-27.png}
\end{figure}

\subsection{\lr{import} کردن کل پکیج}

شما می‌توانید تمام کلاس‌های موجود در یک پکیج را یک‌جا \lr{\texttt{import}}
کنید. این کار را با عبارت \lr{\texttt{*}} می‌توانید انجام بدهید. مثلاً کد
زیر تمام کلاس‌های موجود در پکیج \lr{\texttt{edu.ap.animals}} را
\lr{\texttt{import}} می‌کند.

\begin{codebox}{java}
import edu.ap.animals.*;
\end{codebox}

این قابلیت کار ما را خیلی اوقات آسان می‌کند. ولی انجام این کار همیشه هم
مناسب نیست. \lr{\texttt{import}} کردن دقیق کلاس‌ها علاوه بر خوانایی بیشتر
کدمان، کمی هم زمان کامپایلمان را کمتر می‌کند.

دقت کنید که پکیج‌ها خودشان می‌توانند شامل پکیج باشند؛ همان‌طوری که پوشه‌ها
می‌توانند داخل خود پوشه داشته باشند. اما عبارت \lr{\texttt{*}} فقط
\_کلاس‌\_های متعلق به پکیج را \lr{\texttt{import}} می‌کند و
\emph{sub-package‌} ها را نمی‌کند ( \lr{import} کردن \lr{recursive}
نداریم).

مثلاً در نظر بگیرید که پکیج \lr{\texttt{java.awt}} یک کتابخانه‌ی استاندارد
جاوا است که شامل یک \lr{sub-package} مثل \lr{\texttt{java.awt.event}}
است. کلاس \lr{\texttt{Color}} متعلق به پکیج \lr{\texttt{java.awt}} و
کلاس \lr{\texttt{ActionEvent}} متعلق به زیرپکیج
\lr{\texttt{java.awt.event}} است.

حال شما اگر به هر دوی این کلاس‌ها نیاز دارید باید هر کدام را جدا
\lr{\texttt{import}} کنید. کد زیر اشتباه است:

\begin{codebox}{java}
import java.awt.*;
\end{codebox}

دلیل اشتباه بودنش آن است که در نتیجه‌ی این کد کلاس \lr{\texttt{Color}} ،
ایمپورت می‌شود ولی کلاس \lr{\texttt{ActionEvent}} نه.

کد زیر درست است:

\begin{codebox}{java}
import java.awt.Color;
import java.awt.event.ActionEvent;
\end{codebox}

\subsubsection{دو مثال دیگر: مرور خاطرات}

\paragraph{\texttt{java.util.Scanner}}

احتمالا یکی از اولین برنامه‌هایی که در جاوا نوشتید گرفتن ورودی از کاربر
بوده است. مثلأ کد ساده‌ی زیر را ببینید:

\begin{codebox}{java}
import java.util.Scanner;

public class Main {
    public static void main(String[] args) {
        int x;  
        Scanner scanner = new Scanner(System.in);  
        x = scanner.nextInt();  
    }  
}
\end{codebox}

شاید قبلاً برایتان سوال شده باشد که آن \lr{\texttt{import}} در خط اول چه
کار می‌کند. اکنون احتمالاً می‌توانیم به راحتی به این سوال جواب بدهیم. در
واقع کلاس \lr{\texttt{Scanner}} متعلق به پکیج \lr{\texttt{java.util}}
است و ما چون می‌خواهیم از این کلاس در کدمان استفاده کنیم، آن را
\lr{\texttt{import}} کردیم.

حتی اگر وارد سورس کد این کلاس شوید، عبارت زیر را در خطوط ابتدایی
می‌بینید:

\begin{codebox}{java}
package java.util;
\end{codebox}

\paragraph{\texttt{javax.swing.JFrame}}

در صورتی که با گرافیک کار کرده باشید، می‌دانید وقتی می‌خواستیم یک
\lr{\texttt{frame}} ایجاد کنیم، از کد زیر استفاده می‌کردیم:

\begin{codebox}{java}
import javax.swing.JFrame;

public class Main {
    public static void main(String[] args) {  
        JFrame jFrame = new JFrame();  
    }  
}
\end{codebox}

همان‌طور که می‌توانیم حدس بزنیم، کلاس \lr{\texttt{JFrame}} متعلق به پکیج
\lr{\texttt{javax.swing}} است و چون ما می‌خواهیم از این کلاس در کدمان
استفاده کنیم، آن را \lr{\texttt{import}} کردیم.

می‌توانیم خیلی راحت این را بررسی کنیم. هر موقع خواستید کد یک کلاس را
ببینید (سورس کد جاوا)، می‌توانید روی اسم آن کلاس راست کلیک کنید و از نوار
\lr{Go To} قسمت \lr{Declaration or Usages} را انتخاب کنید. می‌بینید که به
سورس کد آن کلاس منتقل خواهید شد.

\begin{figure}[H]
\centering
\includegraphics[width=0.8\textwidth]{images/oop-1-28.png}
\end{figure}

اگر این کار را انجام دهید، در فایل \lr{\texttt{JFrame.java}} که به آن
منتقل شده‌اید در خطوط ابتدایی عبارت \lr{\texttt{package javax.swing;}} را
می‌بینید.

\subsection{چند نکته در مورد پکیج‌ها}

\subsubsection{هر کلاس دقیقاً به یک پکیج تعلق دارد}

هر کلاس دقیقاً به یک پکیج تعلق دارد. اگر کلاسی که نوشته‌اید را در یک پکیج
قرار ندادید، به صورت پیش‌فرض این کلاس متعلق به \lr{default package} خواهد
بود. مثلاً کلاس زیر را در نظر بگیرید:

\begin{codebox}{java}
public class Student {
    private String name;
    private int age;
    
    // Some other variables and methods
}
\end{codebox}

دقت کنید که اکنون کلاس \lr{\texttt{Student}} را نمی‌شود در پروژه‌ها یا
کلاس‌های دیگر \lr{\texttt{import}} و استفاده کرد؛ چون به پکیج نام‌گذاری
شده‌ای تعلق ندارد ( \lr{default package} واقعاً اسم یک پکیج نیست).

تعریف نکردن پکیج برای کلاس‌هایمان در پروژه‌های کوچک ایرادی ندارد؛ ولی در
پروژه‌های بزرگ باید پکیج‌های مناسبی ایجاد کنیم.

همچنین یک کلاس نمی‌تواند به بیش از یک پکیج تعلق داشته باشد (در غیر این
صورت خطای کامپایل دریافت می‌کنیم).

\subsubsection{استفاده از یک کلاس بدون \lr{import} کردن آن}

گاهی اوقات که فقط یک بار می‌خواهید از یک کلاس در کدتان استفاده کنید،
می‌توانید آن کلاس را به طور مستقیم \lr{import} نکنید و به جای آن به طور
کامل به اسم پکیج در کد اشاره کنید. مثال زیر را ببینید:

\begin{codebox}{java}
public class Main {
    public static void main(String[] args) {  
        javax.swing.JFrame jFrame = new javax.swing.JFrame();  
    }  
}
\end{codebox}

در اینجا ما بدون استفاده از عبارت
\lr{\texttt{import javax.swing.JFrame;}} توانستیم از کلاس
\lr{\texttt{JFrame}} استفاده کنیم. این‌گونه نوشتن شاید کمی طولانی به نظر
بیاید، ولی در کلاس‌هایی که تعداد زیادی \lr{\texttt{import}} دارند و یا
قرار باشد کلاس‌هایی با نام یکسان \lr{\texttt{import}} بشوند، به ما کمک
می‌کند این مشکلات را حل کنیم.

\section{چه چیزی یاد گرفتیم؟}

در این داکیومنت، با مفاهیم پایه‌ای شی‌گرایی در جاوا آشنا شدیم و قدم اول
برای نوشتن برنامه‌های ساخت‌یافته‌تر را برداشتیم. مهم‌ترین نکاتی که یاد
گرفتیم:

\begin{itemize}
\tightlist
\item
  تفاوت بین \textbf{\lr{primitive type}} و \textbf{\lr{reference type}}
  و نحوه‌ی نگه‌داری آن‌ها در حافظه
\item
  این که چگونه می‌توان با استفاده از \lr{\texttt{class}} ،
  \textbf{\lr{type} های جدید} تعریف کرد
\item
  مفهوم \textbf{\lr{object (instance)}} و نحوهٔ ساخت آن با استفاده از
  \lr{\texttt{new}}
\item
  آشنایی با \textbf{\lr{field}} ها و نقش آن‌ها در نگه‌داری وضعیت یک شی
\item
  تعریف و استفاده از \textbf{\lr{method}} ها برای پیاده‌سازی رفتار کلاس‌ها
\item
  نوشتن و استفاده از \textbf{\lr{constructor}} ها برای مقداردهی اولیه‌ی
  \lr{object} ها
\item
  تفاوت بین \textbf{متغیرهای محلی} و \textbf{\lr{field} ها}
\item
  مفهوم \textbf{\lr{static}} و تفاوت اعضای استاتیک با اعضای معمولی
\item
  آشنایی با \textbf{\lr{this}} و کاربرد آن در جلوگیری از \lr{shadowing}
\item
  مفهوم \textbf{\lr{method overloading}} و استفاده از چند متد هم‌نام با
  ورودی‌های متفاوت
\item
  درک بهتر از نحوه‌ی کار \textbf{\lr{reference} ها} و تاثیر آن‌ها در تغییر
  داده‌ها
\item
  آشنایی مقدماتی با \textbf{\lr{Garbage Collection}} و مدیریت خودکار
  حافظه در جاوا
\item
  اهمیت \textbf{\lr{package}}ها برای سازمان‌دهی کد و استفاده از
  \lr{\texttt{import}}
\end{itemize}

در ادامه‌ی مسیر، روی همین مفاهیم پایه، مباحث مهم‌تری مثل
\textbf{کپسوله‌سازی \lr{(Encapsulation)}} ، \textbf{وراثت
\lr{(Inheritance)}} و \textbf{چندریختی \lr{(Polymorphism)}} را یاد
خواهیم گرفت.
	
	
	
	
%%%%%%%%%%%%%%%%%%%%%%%%%%%%%%%%%%%%%%%%%%%%%%%%%%%%%%%%%%%%
\end{document}